\documentclass[12pt,a4paper]{article}
\usepackage[a4paper,margin=18mm]{geometry}
\usepackage[T2A]{fontenc}
\usepackage[utf8]{inputenc}
\usepackage[russian]{babel}
\usepackage{amsmath,amssymb,mathtools}
\usepackage{array,booktabs,longtable}
\usepackage{xcolor}
\usepackage{hyperref}
\usepackage{tikz}

\hypersetup{
  colorlinks=true,
  linkcolor=blue!55!black,
  urlcolor=blue!55!black
}
\setlength{\parindent}{0pt}
\setlength{\parskip}{6pt}
\renewcommand{\arraystretch}{1.25}

\begin{document}

\begin{titlepage}
\centering
\vspace*{0.24\textheight}

{\LARGE\bfseries Ревью домашней работы\par}
\vspace{18mm}
{\large Автор: Лёвин Артём Александрович\par}
\vspace{7mm}
{\large 12 сентября 2026 года\par}

\vfill
\begin{tikzpicture}[scale=1]
  \draw[blue!55!black, line width=0.8pt]
    (0,0) .. controls (2,0.45) and (4,-0.45) .. (6,0)
          .. controls (8,0.45) and (10,-0.45) .. (12,0);
\end{tikzpicture}
\vspace*{8mm}
\end{titlepage}

\newpage
\section*{Сводка}
\small
\begin{tabular}{@{}p{0.08\linewidth}p{0.18\linewidth}p{0.28\linewidth}p{0.16\linewidth}p{0.16\linewidth}@{}}
\toprule
\textbf{№} & \textbf{Тема} & \textbf{Суть ошибки} & \textbf{Ваш ответ} & \textbf{Правильный ответ} \\
\midrule
\hyperref[task:3]{3} &
Показательная функция; граничное значение &
Решение остановлено после записи граничного равенства: отношение \(12{,}5/100\) не приведено к степени двойки, поэтому граница времени и итоговый ответ не получены. &
не указан &
\(15\) единиц времени \\
\bottomrule
\end{tabular}
\normalsize

\newpage
\thispagestyle{empty}
\vspace*{\fill}
\begin{center}
{\LARGE\bfseries Разбор ошибок}
\end{center}
\vspace*{\fill}

\newpage
\phantomsection\label{task:3}
\section*{Задание 3}

\subsection*{Текст задания}
Масса вещества меняется по закону
\[
M(t)=100\cdot 2^{-t/5}.
\]
В течение скольких единиц времени масса будет не меньше \(12{,}5\)?

\subsection*{Ваше решение}
В тетради верно записан закон изменения массы:
\[
M(t)=100\cdot 2^{-t/5},
\]
а затем составлено граничное равенство
\[
12{,}5=100\cdot 2^{-t/5}.
\]
После этого есть несколько зачёркнутых записей с дробями. Не совсем понятно, что именно было записано в этих строках, однако это не мешает проверке: решение на этом этапе не доведено до значения \(t\).

\subsection*{Ваш ответ}
Не указан.

\subsection*{Ошибка}
\textcolor{red}{Ошибка:} решение не доведено после правильной записи граничного равенства.

\subsection*{Где именно возникает ошибка}
Последний полностью корректный шаг:
\[
12{,}5=100\cdot 2^{-t/5}.
\]
Следующий необходимый шаг должен был состоять в делении обеих частей на \(100\) и приведении полученной дроби к степени двойки:
\[
2^{-t/5}=\frac{12{,}5}{100}
=\frac{125}{1000}
=\frac18
=2^{-3}.
\]
Именно этот переход в решении не завершён, поэтому значение \(t\) не найдено.

\subsection*{Почему это важно}
Граничное равенство выбрано правильно: оно позволяет найти момент, когда масса ровно равна \(12{,}5\). Но само равенство ещё не даёт численного ответа.

Здесь нужно заметить, что
\[
\frac{12{,}5}{100}=0{,}125=\frac18,
\]
а
\[
\frac18=2^{-3}.
\]
После этого обе части равенства имеют одинаковое основание \(2\). Так как показательная функция с основанием \(2\) принимает одинаковые значения только при одинаковых показателях, можно приравнять показатели:
\[
-\frac{t}{5}=-3.
\]

Кроме того, в условии спрашивается не только момент равенства, а весь промежуток времени, когда масса \emph{не меньше} \(12{,}5\). Поэтому после нахождения границы нужно ещё определить, по какую сторону от неё выполняется условие. При увеличении \(t\) показатель \(-t/5\) уменьшается, следовательно, \(2^{-t/5}\) и вся масса \(M(t)\) уменьшаются. Значит, до граничного момента масса больше либо равна \(12{,}5\), а после него становится меньше \(12{,}5\).

\subsection*{Ключевая идея}
\textcolor{blue!60!black}{Ключевая идея:} представить отношение \(12{,}5/100\) как степень двойки. Число \(1/8\) равно \(2^{-3}\), поэтому после деления на \(100\) можно сразу сравнить показатели степени.

\subsection*{Исправление от последнего правильного шага}
Продолжаем с уже правильно записанного равенства:
\[
12{,}5=100\cdot 2^{-t/5}.
\]
Делим обе части на \(100\):
\[
2^{-t/5}=\frac{12{,}5}{100}.
\]
Преобразуем правую часть:
\[
\frac{12{,}5}{100}=\frac18=2^{-3}.
\]
Получаем:
\[
2^{-t/5}=2^{-3}.
\]
Основания одинаковы, поэтому показатели равны:
\[
-\frac{t}{5}=-3.
\]
Умножаем обе части на \(-5\):
\[
t=15.
\]
Это граничный момент, когда масса становится равной \(12{,}5\).

\subsection*{Полное правильное решение}
Нужно определить, когда
\[
M(t)\ge 12{,}5.
\]
Подставляем формулу:
\[
100\cdot 2^{-t/5}\ge 12{,}5.
\]
Делим обе части на положительное число \(100\), поэтому знак неравенства не меняется:
\[
2^{-t/5}\ge \frac{12{,}5}{100}.
\]
Правую часть представляем как степень двойки:
\[
\frac{12{,}5}{100}=\frac18=2^{-3}.
\]
Тогда
\[
2^{-t/5}\ge 2^{-3}.
\]
Основание \(2>1\), поэтому функция \(2^x\) возрастает и можно сравнить показатели:
\[
-\frac{t}{5}\ge -3.
\]
Умножаем обе части неравенства на отрицательное число \(-5\). При этом знак неравенства меняется:
\[
t\le 15.
\]
Так как \(t\) обозначает время, рассматриваем \(t\ge 0\). Следовательно,
\[
0\le t\le 15.
\]
Значит, условие выполняется в течение первых \(15\) единиц времени.

\subsection*{Проверка}
При \(t=15\):
\[
M(15)=100\cdot 2^{-15/5}
=100\cdot 2^{-3}
=100\cdot\frac18
=12{,}5.
\]
Граничное значение найдено верно.

Если \(0\le t<15\), то
\[
-\frac{t}{5}>-3,
\]
поэтому
\[
2^{-t/5}>2^{-3},
\]
и, следовательно,
\[
M(t)>12{,}5.
\]
Если \(t>15\), то показатель становится меньше \(-3\), поэтому масса становится меньше \(12{,}5\). Значит, найденный промежуток действительно полный.

\subsection*{Итог}
\textcolor{teal!60!black}{Итог:} масса будет не меньше \(12{,}5\) при
\[
0\le t\le 15,
\]
то есть в течение \(15\) единиц времени.

\subsection*{Задача на закрепление}
Масса вещества меняется по закону
\[
M(t)=160\cdot 2^{-t/4}.
\]
В течение скольких единиц времени масса будет не меньше \(20\)?

\subsection*{Краткое сообщение}
Анастасия, в задании 3 начало решения выбрано верно: Вы правильно записали граничное равенство. Дальше не хватило одного ключевого преобразования: отношение двенадцати с половиной к ста нужно увидеть как одну восьмую, то есть как двойку в степени минус три. После этого показатели легко сравниваются, и получается граница в пятнадцать единиц времени. Стоит отдельно потренировать именно переход от десятичных дробей к удобной степени основания.

\vfill
\begin{center}
\small Лёвин Артём Александрович эксклюзивно для Анастасии Павловой
\end{center}

\end{document}
